\documentclass[12pt]{article}
\usepackage[utf8]{inputenc}
\usepackage[T1]{fontenc}
\usepackage{hyperref}
\title{Morrobalama Language Resources for K-12 Curriculum Development}
\author{Compiled by Claude Sonnet 5 (claude-sonnet-5) for Andr\'as Kornai}
\date{}
\begin{document}
\maketitle

\section{Introduction}
Morrobalama (also Umbuygamu or Morrobolam; ISO 639-3 umg) is a Pama-Nyungan language from the east coast of Cape York Peninsula, Queensland. It was spoken around Princess Charlotte Bay. By the 1990s only four elderly speakers remained. Search results mention community revival efforts supported by the Pama Language Centre.

\section{Primary Educational Resources}
The search found no primers, textbooks or school programmes.

\section{Linguistic and Reference Materials}
\begin{itemize}
\item Sarah Ogilvie (1994), \emph{The Morrobalama (Umbuygamu) language of Cape York Peninsula, Australia}, ANU thesis, the main descriptive source: \url{https://openresearch-repository.anu.edu.au/handle/1885/110346}
\item Wikipedia article ``Morrobolam language'': \url{https://en.wikipedia.org/wiki/Morrobolam_language}
\item Journal of the International Phonetic Association, ``Umbuygamu'' (Illustration of the IPA): \url{https://resolve.cambridge.org/core/journals/journal-of-the-international-phonetic-association/article/umbuygamu/47A9FF39BEF373DC768F0140FBDFCFD1}
\end{itemize}

\section{Religious and Community Resources}
None found in the search.

\section{Digital and Audio Resources}
State Library of Queensland language blog entry on Umbuygamu: \url{https://www.slq.qld.gov.au/blog/language-week-week-thirty-umbuygamu}

\section{Curriculum Development Context}
Descriptive documentation exists, notably the Ogilvie thesis. No teaching materials were found. Speakers were displaced in the early 1960s. A revival would depend on descendant communities and the Pama Language Centre. I did not find a contact for either.
\end{document}